\documentclass[11pt]{article}
\usepackage{mathblatt}
\usepackage{multicol}


\begin{document}
\pfheftstil
\pfheftkopf{Grundwert $G$ \textperiodcentered{} Lösungen}{}
\begin{pfloesung}{}
\lz{1.}{1 \% = $\tfrac{1}{100}$ = 0,01; 10 \% = $\tfrac{1}{10}$ = 0,1; 20 \% = $\tfrac{1}{5}$ = 0,2; 25 \% = $\tfrac{1}{4}$ = 0,25; 50 \% = $\tfrac{1}{2}$ = 0,5; 100 \% = 1 = 1}{}{}
\end{pfloesung}
\begin{pfloesung}{}
\lz{2.}{7 €}{}{}
\end{pfloesung}
\begin{pfloesung}{}
\lz{3.}{$800\,€$}{$1\,\% = 8\,€$; $100\,\%$}{}
\end{pfloesung}
\begin{pfloesung}{}
\lz{4.}{$160$ m, $80$ m, $64$ m}{Ganzes bei $10\,\%$: $16 \cdot 100 : 10 = 160$; Ganzes bei $20\,\%$: $16 \cdot 100 : 20 = 80$; Ganzes bei $25\,\%$: $16 \cdot 100 : 25 = 64$}{}
\end{pfloesung}
\begin{pfloesung}{}
\lz{5.}{800 €}{G = W : p = 200 € : 0,25}{}
\end{pfloesung}
\begin{pfloesung}{}
\lz{6.}{30 €}{G = W : p = 6 € : 0,2}{}
\end{pfloesung}
\begin{pfloesung}{}
\lz{7.}{450 € (360 € sind 80 \%)}{}{}
\end{pfloesung}
\begin{pfloesung}{}
\lz{8.}{G = 4 : $\tfrac{2}{3}$ = 6 Kugeln ($\tfrac{1}{3}$ sind 2 Kugeln)}{}{}
\end{pfloesung}
\begin{pfloesung}{}
\lz{9.}{G = 8 cm : 0,93 $\approx$ 8,60 cm}{}{}
\end{pfloesung}
\end{document}